% =====================================================================
% 6. Análisis exploratorio
% =====================================================================
\section{Análisis exploratorio de las huellas}
\label{sec:eda}

El EDA se hizo sobre los dos Parquet completos, sin muestreo
(\ruta{notebooks/eda_edificios.ipynb}); las figuras se regeneran con
\campo{src/analisis/figuras\_eda.py}. Sus cifras coinciden con los documentos
cargados en MongoDB.

\subsection{Totales y tamaño de las huellas}

\begin{table}[H]
\centering\small
\caption{Totales por fuente sobre los 170 municipios PDET (áreas en EPSG:9377).}
\label{tab:eda-totales}
\begin{tabularx}{\textwidth}{@{}Lrrrr@{}}
\toprule
\textbf{Fuente} & \textbf{Edificios} & \textbf{Área total (km\textsuperscript{2})} & \textbf{Área media (m\textsuperscript{2})} & \textbf{Mediana (m\textsuperscript{2})} \\
\midrule
Google Open Buildings & 2\,694\,786 & 222,17 & 82,45 & 57,87 \\
Microsoft Building Footprints & 1\,761\,048 & 221,13 & 125,57 & 76,99 \\
\bottomrule
\end{tabularx}
\end{table}

Google tiene 1,53 veces más edificios que Microsoft, pero más pequeños: su área
media es un tercio menor y su mediana un 25\,\% menor. Aun así, el área total
es casi igual (diferencia de 1,04\,km\textsuperscript{2}, menos del 0,5\,\%):
Google reparte la misma superficie en más huellas. En Microsoft la cola de
edificios grandes es más larga (p99 de 912\,m\textsuperscript{2} frente a
445\,m\textsuperscript{2}).

\fig{eda_confianza_area.png}{0.28}{Izquierda: distribución de áreas por fuente
(escala logarítmica, medianas marcadas). Derecha: distribución de la confianza
de Google, con los umbrales 0,70 y 0,75 señalados.}{fig:eda-confianza}

La confianza de Google va de 0,65 a 0,98 (media y mediana 0,78). Ningún
registro queda por debajo de 0,65, porque Google solo publica detecciones por
encima de ese valor; por eso 0,65 es el escenario base y 0,70 y 0,75 se usan
para medir la sensibilidad (sección~\ref{sec:sensibilidad}).

\subsection{Validación del área calculada}

Para Google se comparó \campo{area\_m2} (calculada en EPSG:9377) con
\campo{area\_m2\_fuente} (la de Google). La correlación es prácticamente
perfecta y la diferencia relativa tiene mediana 0 y percentiles 5 y 95 de
$\pm$1\,\%; las diferencias mayores son huellas muy pequeñas, donde un redondeo
pesa mucho. El cálculo de área no tiene sesgos sistemáticos.

\fig{eda_comparacion.png}{0.26}{Izquierda: área calculada frente a la de Google
(recta 1:1). Centro: diferencia relativa. Derecha: área total por fuente.}{fig:eda-comparacion}

\subsection{Distribución por subregión y cobertura}

Las dos fuentes concentran sus edificios en las mismas subregiones, pero en
distinta proporción: la Sierra Nevada - Perijá tiene 645\,009 edificios en
Google y 295\,383 en Microsoft. El Chocó es la subregión con menos edificios en
Google (20\,722) y no en Microsoft (37\,692): la primera señal de que las
fuentes no cubren el territorio igual.

\fig{eda_subregion.png}{0.3}{Edificios por subregión PDET y fuente.}{fig:eda-subregion}

El mapa de densidad de centroides (figura~\ref{fig:eda-mapa}) muestra dónde
están los huecos: cinco municipios no tienen ningún edificio en Google
(Murindó y Vigía del Fuerte en Antioquia; Medio San Juan, Nóvita y Bojayá en el
Chocó), mientras que Microsoft tiene entre 1\,331 y 2\,004 en cada uno. Otros
tienen muy pocos en Google, como Sipí (1), López de Micay (24) y Condoto (239).

\fig{eda_mapa_densidad.png}{0.4}{Densidad de centroides por fuente sobre los
municipios PDET (escala logarítmica común). En rojo, los cinco municipios sin
edificios de Google.}{fig:eda-mapa}

Las figuras de esta sección son la versión corregida de las del EDA original,
rehechas con un estilo común para todo el informe (textos en español, escalas
logarítmicas reales, miles separados y la misma escala de color en los dos
paneles), atendiendo la retroalimentación recibida sobre las gráficas.
